\sect{द्वादशी-एकादशी-जागरणमहिमावर्णनम्}
\label{sec:padma-jagaranamahima}

\ganapatyadiDhyanam
\padmaPuranam

\dnsub{अथ नामैकोनचत्वारिंशोऽध्यायः॥३९}

\uvacha{महादेव उवाच}

\twolineshloka
{शृणु नारद वक्ष्यामि माहात्म्यं जागरस्य च}
{यच्छ्रुत्वा मुक्तिमाप्नोति महापापी न संशयः}% ॥१॥

\uvacha{नारद उवाच}

\twolineshloka
{अहो विश्वेश्वरो विष्णुः पवित्रीकरणः सदा}
{तस्योपवासमाहात्म्यं श्रुतं हि त्वन्मुखाच्छिव}% ॥२॥

\twolineshloka
{तथाऽपि श्रोतुमिच्छामि माहात्म्यं जागरस्य तु }
{कीदृक्जागरमाहात्म्यं रात्रो भक्तिस्तु  कीदृशी}% ॥३॥

\threelineshloka
{प्रहरेषु च या पूजा वद विश्वेश्वर प्रभो}
{त्वं लोकेषु सदा पूज्यस्त्वं हि देवो जनार्दनः}
{त्वं हि विश्वेश्वरो देवो यतो भक्तिर्जनार्दने}% ॥४॥

\twolineshloka
{सर्वेषां चैव भक्तानां त्वं च श्रेष्ठ उमापतिः}
{लोकेस्मिन् सर्वदा भक्त्या तवाख्या वर्तते सदा}% ॥५॥

\twolineshloka
{अतो येन प्रकारेण लोकानां मुक्तिरेव च}
{विश्वेश्वर वद त्वं तु माहात्म्यं जागरस्य तु}% ॥६॥

\uvacha{महादेव उवाच}

\twolineshloka
{एकादश्यां जनो विष्णुं रात्रौ सम्पूज्य भक्तितः}
{कुर्याज्जागरणं विष्णोः पुरतो वैष्णवैः सह}% ॥७॥

\twolineshloka
{गीतं वाद्यं तथा नृत्यं पुराणपठनं तथा}
{धूपं दीपं च नैवेद्यं पुष्पं गन्धानुलेपनम्}% ॥८॥

\twolineshloka
{फलमर्घं तथा श्रद्धा दानमिन्द्रियसंयमम्}
{सत्यान्वितं च विप्रेन्द्र वचोयुक्तं क्रियान्वितम्}% ॥९॥

\twolineshloka
{विनिद्रं च मुदायुक्तो यः करोति नरः सदा}
{सर्वपापविनिर्मुक्तो जायते विष्णुवल्लभः}% ॥१०॥

\twolineshloka
{रात्रौ जागरणे प्राप्ते निद्रां कुर्वन्ति वैष्णवाः}
{हारितं चोपवासं तैर्व्रतं वै विष्णुसंज्ञकम्}% ॥११॥

\twolineshloka
{ये कुर्वन्ति नराः प्राज्ञ जागरे विष्णुसंज्ञके}
{जागरं कृष्णभावेन न स्वपन्ति कदाचन}% ॥१२॥

\twolineshloka
{कृष्णस्य नाम मनसा वदन्ति च पुनः पुनः}
{ते तु धन्यतमा ज्ञेया अस्यां रात्रौ विशेषतः}% ॥१३॥

\twolineshloka
{क्षणेक्षणे तु गोदानं घट्यां चैव चतुर्गुणम्}
{प्रहरे कोटिगुणितं चतुर्यामेष्वसङ्ख्यकम्}% ॥१४॥

\twolineshloka
{जागरे निमिषार्धे तु केशवाग्रे विशेषतः}
{तत्फलं कोटिगुणितं तस्य सङ्ख्या न विद्यते}% ॥१५॥

\twolineshloka
{नर्तनं कुरुते यस्तु केशवाग्रे नरोत्तमः}
{न फलं हीयते तस्य आजन्ममरणान्तिकम्}% ॥१६॥

\twolineshloka
{साश्चर्यं चैव सोत्साहं पापालापादिवर्जितम्}
{प्रदक्षिणसमायुक्तं नमस्कारपुरःसरम्}% ॥१७॥

\twolineshloka
{नीराजनसमायुक्तमनिर्विण्णेन चेतसा}
{यामेयामे महाभाग कुर्यादारार्तिकं हरेः}% ॥१८॥

\twolineshloka
{षड्विंशगुणसंयुक्तमेकादश्यां च जागरम्}
{यः करोति नरो भक्त्या न पुनर्जायते भुवि}% ॥१९॥

\twolineshloka
{य एवं कुरुते भक्त्या वित्तशाठ्यविवर्जितः}
{जागरं वासरे विष्णोर्लीयते परमात्मनि}% ॥२०॥

\twolineshloka
{धनवान्वित्तशाठ्येन यः करोति प्रजागरम्}
{तेनात्माहारितो नूनं कितवेन दुरात्मना}% ॥२१॥

\twolineshloka
{विष्णुजागरणे प्राप्ते उपहासं करोति यः}
{षष्टिवर्षसहस्राणि विष्ठायां जायते कृमिः}% ॥२२॥

\twolineshloka
{वेदविद् ब्राह्मणो यस्तु नर्तनेन विशेषतः}
{उपहासपरः प्राप्तः स वै चाण्डाल उच्यते}% ॥२३॥

\twolineshloka
{निमिषं निमिषार्धं वा यः करोति प्रजागरम्}
{धर्मार्थकाममोक्षाणां प्राप्नोति पदमव्ययम्}% ॥२४॥

\twolineshloka
{वेदशास्त्ररतो नित्यं नित्यं वै यज्ञयाजकः}
{रात्रौ जागरणे प्राप्ते निन्दां कुर्वन्व्रजत्यधः}% ॥२५॥

\twolineshloka
{मम पूजां प्रकुर्वाणो विष्णुनिन्दासु तत्परः}
{एकविंशत्कुलेनैव नरकं प्रतिपद्यते}% ॥२६॥

\threelineshloka
{विष्णुः शिवः शिवो विष्णुरेकमूर्तिर्द्विधा स्थिताः}
{तस्मात् सर्वप्रकारेण नैव निन्दां प्रकारयेत्}
{दष्टाः कलिभुजङ्गेन स्वपन्ति मधुहाहनि}% ॥२७॥

\twolineshloka
{कुर्वन्ति जागरं नैव मायया ते च मोहिताः}
{प्राप्ता एकादशी येषां कलौ जागरणं विना}% ॥२८॥

\twolineshloka
{ते विनष्टा न सन्देहो यस्माज्जीवितमध्रुवम्}
{उद्धृतं नेत्रयुग्मं तु दत्त्वा वै वैष्णवं पदम्}% ॥२९॥

\twolineshloka
{कृतं ये नैव पश्यन्ति पापिनो हरिजागरम्}
{अभावे वाचकस्याथ गीतं नृत्यं तु कारयेत्}% ॥३०॥

\twolineshloka
{वाचके सति देवर्षे पुराणं प्रथमं पठेत्}
{अश्वमेधसहस्रस्य वाजपेयायुतस्य च}% ॥३१॥

\twolineshloka
{पुण्यं कोटिगुणं वत्स विष्णोर्जागरणे कृते}
{पितृपक्षे मातृपक्षे भार्यापक्षे तु वाडव}% ॥३२॥

\twolineshloka
{कुलानुद्धरते चैतान्कृत्वा जागरणं हरेः}
{उपोषणदिने विद्धे जागरं पूजनं हरेः}% ॥३३॥

\twolineshloka
{वृथादानादिकं सर्वं कृतघ्नेषु कृतं यथा}
{उपोषणदिने विद्धे प्रारब्धे जागरे स्थितिम्}% ॥३४॥

\twolineshloka
{विहाय स्थानं तद् विष्णुः शापं दत्वा प्रगच्छति}
{अविद्धे वासरे विष्णोर्ये कुर्वन्ति प्रजागरम्}% ॥३५॥

\twolineshloka
{तेषां मध्ये तु तुष्टः सन्नृत्यं तु कुरुते हरिः}
{यावद्दिनानि कुरुते जागरं केशवाग्रतः}% ॥३६॥

\twolineshloka
{युगानि तानि तावन्ति विष्णुलोके महीयते}
{यावद्दिनानि वसते विना जागरणं हरेः}% ॥३७॥

\twolineshloka
{तावद् वर्षसहस्राणि रौरवान्न निवर्तते}
{एकदश्यां शयानस्तु विना जागरणं हरेः}% ॥३८॥

\twolineshloka
{मूकवत्तिष्ठते यो वै गानं पाठं न वाचरेत्}
{सप्तजन्मानि मूकत्वं जायतेऽजागरे हरेः}% ॥३९॥

\twolineshloka
{यो न नृत्यति मूढात्मा पुरतो जागरे हरेः}
{पङ्गुत्वं तस्य जानीयात्सप्तजन्मानि वाडव}% ॥४०॥

\twolineshloka
{यः पुनः कुरुते गीतं नृत्यं जागरणं हरेः}
{ब्राह्मं पदं मदीयं च सत्यं वै तस्य वैष्णवम्}% ॥४१॥

\twolineshloka
{यः प्रबोधयते लोकान्विष्णोर्जागरणे रतः}
{वसेच्चिरं तु वैकुण्ठे पितृभिः सह वैष्णवः}% ॥४२॥

\twolineshloka
{मतिं प्रयच्छते यस्तु हरेर्जागरणं प्रति}
{षष्टिवर्षसहस्राणि श्वेतद्वीपे वसेन्नरः}% ॥४३॥

\twolineshloka
{यत्किञ्चित्क्रियते पापं कोटिजन्मनि मानवैः}
{श्रीकृष्णजागरे सर्वं रात्रौ नश्यति वाडव}% ॥४४॥

\twolineshloka
{शालग्रामशिलाग्रे ये कुर्वन्ति प्रतिजागरम्}
{यामेयामे फलं प्रोक्तं कोट्यैन्दवसमुद्भवम्}% ॥४५॥

\twolineshloka
{सम्प्राप्ते वासरे विष्णोर्ये न कुर्वन्ति जागरम्}
{वृथा स्यात्तत्कृतं तेषां वैष्णवानां च निन्दया}% ॥४६॥

\twolineshloka
{कामार्थौ सम्पदः पुत्राः कीर्तिर्लोकाश्च शाश्वताः}
{यज्ञायुतैर्न लभ्यन्ते द्वादशीजागरं विना}% ॥४७॥

\twolineshloka
{मतिर्न जायते यस्य द्वादश्यां जागरं प्रति}
{न हि तस्याधिकारोऽस्ति पूजने केशवस्य हि}% ॥४८॥

\twolineshloka
{यावत् पदानि चलति केशवायतनं प्रति}
{अश्वमेधसमानि स्युः जागरार्थं प्रगच्छतः}% ॥४९॥

\twolineshloka
{पादयोः पतितं यावद्धरण्यां पांशु गच्छताम्}
{तावद् वर्षसहस्राणि जागरो वसते दिवि}% ॥५०॥

\twolineshloka
{तस्माद् गृहात् प्रगन्तव्यं जागरे केशवालये}
{कलौ मलविनाशाय द्वादशीद्वादशीषु च}% ॥५१॥

\twolineshloka
{परापवादसंयुक्तं मनः प्रासाद वर्जितम्}
{शास्त्रहीनमगान्धर्वं तथा दीपविवर्जितम्}% ॥५२॥

\twolineshloka
{शक्त्योपचाररहितमुदासीनं सनिन्दनम्}
{कलियुक्तं विशेषेण जागरं नवधा मतम्}% ॥५३॥

\twolineshloka
{सशास्त्रं जागरं यच्च नृत्यगान्धर्वसंयुतम्}
{सवाद्यं तालासंयुक्तं सदीपं मधुभिर्युतम्}% ॥५४॥

\twolineshloka
{उच्चारैस्तु समायुक्तं यथोक्तैर्भक्तिभावितैः}
{प्रसन्नं तुष्टिजननं सम्मूढं लोकरञ्जनम्}% ॥५५॥

\twolineshloka
{गुणैर्द्वादशभिर्युक्तं जागरं माधवप्रियम्}
{कर्तव्यं तत् प्रयत्नेन पक्षयोः शुक्लकृष्णयोः}% ॥५६॥

\twolineshloka
{किं व्रतैर्बहुभिश्चीर्णैस्तीर्थवासेन तस्य किम्}
{द्वादशीवासरे प्राप्ते न कुर्याज्जागरं हरेः}% ॥५७॥

\twolineshloka
{प्रवासेन त्यजेद्यस्तु पथिस्विन्नोऽपि वाडव}
{जागरं वासुदेवस्य द्वादश्यां तु समे प्रियः}% ॥५८॥

\twolineshloka
{मद्भक्तो न हरेः कुर्याज्जागरं पापमोहितः}
{व्यर्थं मत्पूजनं तस्य मत्पूज्यं यो न पूजयेत्}% ॥५९॥

\twolineshloka
{न शैवो न च सौरोऽसौ न शाक्तो गणसेवकः}
{यो भुङ्क्ते वासरे विष्णोर्ज्ञेयः पश्वधिको हि सः}% ॥६०॥

\twolineshloka
{विप्रियं च कृतं तेन दुष्टेनैव च पापिना}
{मद्भक्तिबलमाश्रित्य यो भुङ्क्ते वै हरेर्दिने}% ॥६१॥

\twolineshloka
{सबाह्याभ्यन्तरं देहं वेष्टितं पापकोटिभिः}
{मुच्यन्ते वासरे विष्णोर्ये कुर्वन्ति प्रजागरम्}% ॥६२॥

\twolineshloka
{कूर्परं यमदूतानां दत्तं तेन यमस्य च}
{कृत्वा जागरणं विष्णोरविद्धं द्वादशीव्रतम्}% ॥६३॥

\twolineshloka
{स्वर्गापेक्षा मुनिश्रेष्ठ मुक्ता ते नैव संशयः}
{वाञ्छितं नारकं सौख्यं विद्धं कृत्वा हरेर्दिनम्}% ॥६४॥

\twolineshloka
{निहताः पितरस्तेन देवानां वै वधः कृतः}
{दत्तं राज्यं तु दैत्यानां कृत्वा विद्धं हरेर्दिनम्}% ॥६५॥

\twolineshloka
{यो नृत्यति प्रहृष्टात्मा कृत्वा वै करताडनम्}
{गीतं कुर्वन्मुखेनापि दर्शयन्कौतुकान्बहून्}% ॥६६॥

\twolineshloka
{पुरतो वासुदेवस्य रात्रौ जागरणे स्थितः}
{पठन्कृष्णचरित्राणि रञ्जयन्वैष्णवान्गणान्}% ॥६७॥

\twolineshloka
{मुखेन कुरुते वाद्यं सम्प्रहृष्टतनूरुहः}
{दर्शयन्विविधान्भृत्यान्स्वेच्छालापान् प्रकारयन्}% ॥६८॥

\twolineshloka
{भावैरेतैर्नरो यस्तु कुरुते जागरे हरेः}
{निमिषेनिमिषे पुण्यं तीर्थकोटिफलं स्मृतम्}% ॥६९॥

\twolineshloka
{अनुद्विग्नमना यस्तु धूपनीराजनं हरेः}
{कुरुते जागरे रात्रौ सप्तद्वीपाधिपो भवेत्}% ॥७०॥

\twolineshloka
{यानि कानि च पापानि ब्रह्महत्यासमानि च}
{कृष्णा ह जागरात्तानि विलयं यान्ति खण्डशः}% ॥७१॥

\twolineshloka
{एकतः क्रतवः सर्वे समाप्तवरदक्षिणाः}
{एकतो देवदेवस्य जागरः कृष्णवल्लभः}% ॥७२॥

\twolineshloka
{तत्र काशी पुष्करं च प्रयागं नैमिषं गया}
{शालग्राममहाक्षेत्रमर्बुदारण्यमेव च}% ॥७३॥

\twolineshloka
{पौष्करं मथुरा तत्र सर्वतीर्थानि चैव हि}
{यज्ञा वेदाश्च चत्वारो व्रजन्ति हरिजागरम्}% ॥७४॥

\twolineshloka
{गङ्गा सरस्वती तापी यमुना च शतद्रुका}
{चन्द्रभागा वितस्ता च नद्यः सर्वास्तु तत्र वै}% ॥७५॥

\twolineshloka
{सरांसि च ह्रदाः सर्वे समुद्राः सर्व एव हि}
{एकादश्यां द्विजश्रेष्ठ गच्छन्ति कृष्णजागरम्}% ॥७६॥

\twolineshloka
{स्पृहणीया हि देवानां ये नराः कृष्णजागरे}
{नृत्यन्ति गीतं कुर्वन्ति वीणावाद्यप्रहर्षिताः}% ॥७७॥

\twolineshloka
{एवं जागरणं कृत्वा सम्पूज्य च महाहरिम्}
{द्वादश्यां पारणं कृत्वा स्वशक्त्या वैष्णवैः सह}% ॥७८॥

\uvacha{महादेव उवाच}

\twolineshloka
{शृणु ब्रह्मन् प्रवक्ष्यामि द्वादशीमाहात्म्यमुत्तमम्}
{द्वादशी तु सदा ज्ञेया पुत्रदा मोक्षदायिनी}% ॥७९॥

\twolineshloka
{प्रातः स्नात्वा हरिं पूज्य उपवासं समर्पयेत्}
{अज्ञानतिमिरान्धस्य व्रतेनानेन केशव}% ॥८०॥

\twolineshloka
{प्रसीद सुमुखो भूत्वा ज्ञानदृष्टिप्रदो भव}
{पारणं च ततः कुर्याद् यथासम्भवमग्रतः}% ॥८१॥

\twolineshloka
{अत ऊर्ध्वं यथेष्टं तु कारयेच्च यथाविधि}
{यदा तु द्वादशी स्वल्पा पारणेन भवेद्दिवज}% ॥८२॥

\twolineshloka
{तदा रात्रौ तु कर्तव्यं पारणं मुक्तिमिच्छता}
{तदा न रात्रिदोषः स्यान्निषिद्धं न भवेत्क्वचित्}% ॥८३॥

\twolineshloka
{यदुक्तं निशि न स्नायान्महानिशि न भोजयेत्}
{तत्पूर्वपरयामाभ्यां दिनवत्कर्म कारयेत्}% ॥८४॥

\twolineshloka
{यदा भवति स्वल्पा तु द्वादशी पारणे दिने}
{उषःकाले  द्वयं कुर्यात् प्रातर्मध्याह्निकं तथा}% ॥८५॥

\twolineshloka
{द्वादशी साधिता येन नरेण भुवि सर्वदा}
{तस्य पुण्यमहं वक्तुं न समर्थो विशेषतः}% ॥८६॥

\twolineshloka
{साधयित्वाखिलान्कामान् प्राप्नुयुश्च महाजनाः}
{अम्बरीषादयः सर्वे ये भक्ता भुवि विश्रुताः}% ॥८७॥

\twolineshloka
{द्वादशीं साधयित्वा तु ते गता विष्णुसद्मनि}
{सत्यं सत्यं पुनः सत्यं यदुक्तं तु मया तव}% ॥८८॥

\twolineshloka
{नास्ति विष्णुसमो देवो न तिथिर्द्वादशीसमा}
{अत्र दत्तं च भुक्तं च तथा पूजादिकं च यत्}% ॥८९॥

\twolineshloka
{तत् सर्वं पूर्णतां याति पूजिते माधवे सति}
{किं पुनर्बहुनोक्तेन भक्तानां वल्लभो हरिः}% ॥९०॥

\twolineshloka
{प्रददात्यखिलान्कामान्यावदाभूतसम्प्लवम्}
{द्वादश्यां चैव यद् दत्तं तत् सर्वं सफलं भवेत्}% ॥९१॥

\twolineshloka
{कुरुक्षेत्रेषु यद् दत्तं निष्फलं नैव जायते}
{तद्वच्च द्वादशीदत्तं भवेद् देवर्षिसत्तम}% ॥९२॥

॥इति श्रीपाद्मे महापुराणे पञ्चपञ्चाशत्साहस्र्यां संहितायामुत्तरखण्डे उमापतिनारदसंवादान्तर्गतकृष्णयुधिष्ठिरसंवादे द्वादशी-एकादशी-जागरणमहिमावर्णनं नाम नामैकोनचत्वारिंशोऽध्यायः॥३९॥

\genMangalaShloka

\hyperref[sec:ekadashi_mahatmyam_padma_puranam]{\closesub}
\clearpage

